\documentclass[notitlepage,12pt]{../../cpsc250l}

\title{CPSC 250L Lab 1\\[0.5em] File/Directory Structure and Git Intro\\[0.5em]Name: $\rule{7cm}{0.15mm}$  Section:$\rule{2cm}{0.15mm}$}

\begin{document}

\maketitle

\section{Introduction}

The focus of this lab is to solidify your understanding of file/directory structure, and
introduce you to navigating via the command line interface (CLI) using GitBash.

Follow the directions on the GitLab project webpage, and answer the questions below.

{\bf NOTE: Access https://gitlab.pcs.cnu.edu  NOT gitlab.com or github.com}

\section{Questions}

Answer or initial that you completed as directed.

\begin{enumerate}
\item Part 1 - File/Directory Structure, and Git Bash
	\begin{enumerate}
	\item \initial Initial when you have opened the Windows File Explorer GUI and navigated to the Documents folder.
	\item \initial Initial when you have navigated to the top level folder in File Explorer.
	\item \initial Initial when you have navigated back to the Documents folder in File Explorer.
	\item \initial Initial when you have opened the GitBash Command Line Interface.
	\item \initial Initial when you have changed to the Desktop directory in GitBash.
	\item What is another name for a CLI?  $\rule{4cm}{0.15mm}$
	\item \initial Initial when you have listed files using \texttt{ls}.
	\item \initial Initial when you have listed files using \texttt{ls -altr}.
	\item What is the full path of the Desktop directory?\\ $\rule{8cm}{0.15mm}$
	\item \initial Initial when you have navigated to the root directory in GitBash.
	\item \initial Initial when you have navigated back to the Desktop directory going step by step.
	\item What is the command used to back up one directory from the current directory?\\
	$\rule{8cm}{0.15mm}$
	\item What is the full path of the current directory after backing up one from the Desktop directory?\\ $\rule{8cm}{0.15mm}$
	\item What is the command used to create a new directory in a bash shell?\\
	$\rule{4cm}{0.15mm}$
	\item What is the command used to display the current directory?\\
	$\rule{4cm}{0.15mm}$
	\item What is the command used to go to a user's home directory (denoted with \textasciitilde) ?\\
	$\rule{4cm}{0.15mm}$
	\item \initial Initial when you have created the \texttt{cpsc250L} folder on your Desktop using GitBash.
	\item What is the shortcut that enables us to avoid typing out the full ``cd Desktop'' (or any other file or directory name) every time?\\
	$\rule{6cm}{0.15mm}$
	\end{enumerate}
\vspace{0.5cm}
\item Part 2 - Git Project
	\begin{enumerate}
	\item \initial Initial when you have set your `WCUSER` and `WCPASS` variables for WebCAT on personal GitLab \textbf{group}.
	\item \initial Initial when you have forked \textbf{cpsc250l-lab01} project repo to your \textbf{personal} group.
	\item \initial Initial when you have set your Git credentials.
	\item \initial when you have cloned \textbf{cpsc250l-lab01} project repo into your \texttt{cpsc250L} directory AND verified its contents using \texttt{ls}.
	\item What is the full path of the repo directory that you just cloned?\\ $\rule{10cm}{0.15mm}$
	\item \initial Initial when you have the project open in PyCharm.
	\item \initial Initial when you set the working directory in PyCharm.
	\item \initial Initial when you make the first *push* to GitLab.
	\item \initial Initial when Exercise 1 is complete and pushed
	\item \initial Initial when Exercise 2 is complete and pushed
	\item \initial Initial when you have verified correct on WebCAT!

	\end{enumerate}
\end{enumerate}

If all is initialed and filled out, and correct on WebCAT, then
show your instructor and hand in this form.

You are then free to leave; have a great remainder of the day.
Of course your day should include some additional coding practice!

This lab is worth 100 points:
\begin{itemize}
 \item 35 points completing this worksheet
 \item 35 points WebCAT solution (correctness)
 \item 30 points lab participation (completing during lab or significant progress)
 \item \initial \textbf{Instructor} initial to verify full participation during lab period
\end{itemize}

 Note: Your WebCAT score will only show 35 of 100 until your
 instructor enters the participation score under "Design/Readability" in WebCAT.
 You goal should be 35/35 correctness score on WebCAT (log in and check!),
 and the instructor initial above to indicate worksheet complete during lab.

\end{document}
